% Source: https://artofproblemsolving.com/community/c3h1377519_just_a_quick_common_sense_questions
% Stable identity: aops-topic-1377519-q01
Let each bird eat $k$ worms per minute. Then the number of worms eaten is
\[
\text{worms}=k\cdot\text{birds}\cdot\text{minutes}.
\]
From $4=k\cdot4\cdot4$ we get $k=\dfrac{1}{4}$: each bird eats one worm every $4$ minutes. Then $32$ birds eat $32\cdot\dfrac{1}{4}=8$ worms per minute, and $32$ worms take $\dfrac{32}{8}=\boxed{4}$ minutes.

The proportion is the same rule rearranged: dividing the formula by $k\cdot\text{worms}$ gives
\[
\dfrac{\text{birds}}{\text{worms}}\cdot\text{minutes}=\dfrac{1}{k} ,
\]
a quantity that is the same in every situation with the same eating rate. Writing it for both situations gives $\dfrac{4}{4}\cdot4=\dfrac{32}{32}\cdot y$, so $y=4$. The ``common sense'' answer is right, but the reason is that the number of worms per bird is the same ($1$) in both situations and each bird needs the same time per worm.
